\chapter{Introduction} 

\section{Singularity Formation in Models of 3D Fluid Flows}
The three-dimensional (3D) Navier-Stokes equations are a set of partial differential equations (PDEs) that describe the motion of homogeneous, viscous, incompressible fluids. {\color{red}The Navier-Stokes equations on the domain $\Omega$ are defined as 
\begin{subequations}\label{eq:Navier:Stokes}
\begin{align}
\partial_t \u + \u\cdot \bn \u + \bn p - \nu \bd \u &= 0 \qquad\text{ in } \Omega\times (0,T],\\
\bn \cdot \u &= 0 \qquad\text{ in } \Omega\times [0,T],\\
\u(x,0) &= \be \qquad\text{ in }\Omega
\end{align}
\end{subequations}
\noindent where $\u = [u_1,u_2,u_3]^\intercal$ is the fluid velocity field, $p$ is the scalar pressure, $\nu>0$ is the coefficient of kinematic viscosity, and $T$ is the length of the time interval. The field $\be$ is the initial condition assumed to be divergence-free. Finally, the $\bn = [\partial_1,\partial_2,\partial_3]^T$ operator is used to define the gradient of the pressure, the tensor $[\bn \u]_{ij}=\partial_j u_i$ for $i,j=1,2,3$, and the divergence of the velocity field. Because of their generality}, the Navier-Stokes equations are utilized to model the behaviour of systems in fields as disparate as aero/hydrodynamics, weather prediction, biophysics, and astrophysics \cite{cho_solving_2022}\cite{doering_3d_2009}. Despite the numerous fields that make use of the 3D Navier-Stokes equations, and extensive research examining their properties, {\color{red}it is only known that weak Leray solutions exist globally in time \cite{leray_sur_1934}. However, these solutions may contain singularities and their uniqueness remains unproven in most cases. It has also yet to be proven whether or not classical solutions in the pointwise sense exist globally in time \cite{protas_systematic_2022}. In fact, this is the basis of the well known millennium problem from the Clay Mathematics Institute which offers a substantial monetary prize for the confirmation of the existence, or non-existence, of a classical solution for smooth initial conditions \cite{robinson_navierstokes_2020}}.
It is common for scientific research that focuses the fundamental properties of the Navier-Stokes equations to utilize either the unbounded domain $\Omega := \RRR$, or the periodic 3D torus $\Omega := \TT^3_L$ with $L>0$ \cite{protas_systematic_2022}. {\color{red}For this analysis, we use the 3D torus ($L=1$) with periodic boundary conditions.}
\\~\\
A central question in fluid mechanics is whether or not the Navier-Stokes equations are globally well-posed or if singularities can form from smooth initial conditions. A singularity occurs when the derivative of the velocity field becomes infinite at a single point \cite{moffatt_singularities_2019}. The formation of a singularity in a fluid flow, from smooth initial conditions, would invalidate the Navier-Stokes equations as an accurate description of fluid flows \cite{protas_systematic_2022}. The question of singularity formation from smooth initial conditions {\color{red} also remains open} for the Euler equations \cite{gibbon_three-dimensional_2008}. The Euler equations are the inviscid form of the Navier-Stokes equations, obtained by setting $\nu=0$
\begin{subequations}\label{eq:Euler}
\begin{align}
\partial_t \u + \u\cdot \bn \u + \bn p &= 0 \qquad\text{ in } \Omega\times (0,T],\\
\boldsymbol\nabla \cdot \mathbf{u} &= 0 \qquad\text{ in } \Omega\times [0,T],\\
\u(x,0) &= \be \qquad\text{ in } \Omega.
\end{align}
\end{subequations}
The Euler are often studied because they are simpler to analyse than the Navier-Stokes equations {\color{red} due to the absence of the viscous dissipation}. {\color{red}There are many different methods for determining if singularity formation has occurred, the most well known of which is the Beale-Kato-Majda (BKM) criterion \cite{beale_remarks_1984} which states: A smooth solution $u$ of the Euler system develops a singularity at $t=t_0$ if and only if\\
\begin{equation}\label{eq:beale:kato:majda}
\lim_{t\rightarrow t_0}\int_0^t \|\bds\omega\|_{L^\infty} d\tau = \infty, \qquad \bds\omega = \nabla \times \u,
\end{equation}

\noindent where $\|\cdot\|_{L^\infty}$ is the standard Lebesgue $L^\infty$ norm. A key difference between the Navier-Stokes and Euler equations is that there is more numerical evidence that indicates singularity formation in Euler is possible from smooth initial conditions. 
\\~\\
This includes the results of \cite{brachet_small-scale_1983} and \cite{bustamante_interplay_2012} which presented numerical evidence for singularity formation in the Euler equations on a 3D torus ($L=2\pi$). As well as the results of \cite{hou_potential_2022} which provide numerical evidence for singularity formation in the interior domain of the 3D axisymmetric Euler equations. The breakthrough results of \cite{chen_stable_2023} which proved singularity formation on the boundary of the domain for 3D axisymmetric Euler equations. We make special note of the optimization-based approach of \cite{zhao_systematic_2023}, which utilized a Riemannian conjugate gradient method to identify 'extreme' initial conditions, and produced numerical evidence for singularity formation in Euler flows on the 3D unit torus. Finally, \cite{larios_computational_2018} which provides numerical evidence for singularity formation in the Euler equations based on the Voigt regularization. These final two papers are of special significance because they form the primary basis for the motivation and methodology of this analysis.}

\section{3D Euler-Voigt Equations}

The Euler-Voigt equations are an inviscid regularization of the Euler equations (\ref{eq:Euler}) which have been shown to be globally well-posed and so do not form singularities from smooth initial conditions \cite{cao_global_2006}. The Euler-Voigt equations are defined as
\begin{subequations}\label{eq:Euler:Voigt}
\begin{align}
-\alpha^2\partial_t\bd\ua + \partial_t \ua + \ua\cdot \bn\ua + \bn p_\alpha &= 0 \qquad\text{ in } \Omega\times (0,T],\\ 
\bn \cdot \ua &= 0 \qquad\text{ in } \Omega\times (0,T],\\
\ua(x,0) &= \bea \qquad\text{ in }\Omega,
\end{align}
\end{subequations}
{\color{red}where the regularization is accomplished} with the addition of the term $-\alpha^2\partial_t \bd \ua$ where $\alpha >0$ is the regularization parameter with units of length. This term {\color{red} stabilizes the system without dissipation}, which would remove energy from the system \cite{larios_computational_2018}. {\color{red}The regularization term acts as a spectral filter and suppresses the development of small scale structures with length-scale smaller than $\alpha$ \cite{molfetta_self-truncation_2015} }. We introduce the subscript $\alpha$ to distinguish between solutions of (\ref{eq:Euler}) and (\ref{eq:Euler:Voigt}) {\color{red}i.e., $\u$ and $\ua$}. 
\\~\\
For (\ref{eq:Euler}) the total kinetic energy is conserved so long as the solution remains well-posed in the classical sense
\begin{equation}\label{eq:total:kinetic:energy}
E_k(t) := \frac{1}{2}\|\u\|_{L^2}^2 = E_k(0),
\end{equation} 
where the norm $\|\cdot\|_{L^2}$ here is the standard Lebesgue $L^2$ norm. For (\ref{eq:Euler:Voigt}) the conserved quantity is known to be the $\alpha$-energy. The $\alpha$-energy is defined as
\begin{equation}\label{eq:alpha:Energy}
E_\alpha(t):=\|\ua(t)\|_{L^2}^2 + \alpha^2\|\bn\ua(t)\|_{L^2}^2 = E_\alpha(0).
\end{equation}


\section{Outline of the Thesis}

{\color{red}The goal of this analysis is to search for numerical evidence of singularity formation in the Euler equations from smooth initial conditions. This search will be performed with a PDE-constrained optimization problem which will identify initial condition that produce 'extreme' behaviour in Euler-Voigt flows. In Chapter 2 we present the properties of the Euler-Voigt equations, including the blow-up criteria, scaling property, and manifold in which we search for initial conditions. In Chapter 3 we present the formulation of the optimization problem. This includes the description of our objective functional, the gradient calculation, and the Riemannian conjugate gradient method used to compute our search direction and iterative relation. In Chapter 4 we describe the numerical methods used to perform the calculations, including spatial resolution, time step size, Gevrey parameters and the regularization parameter. In addition we validate the model using the spectrum, kappa test, and conservation of $\alpha-$energy. In Section 5 we present our results. We begin with the results of the Euler-Voigt equations from the Taylor Green Vortex without optimization as $\alpha\rightarrow 0$. These results indicate that singularity formation does not occur with the time window. We then present our optimization results as $\alpha\rightarrow 0$, which indicate that RESULTS RESULTS RESULTS. Finally, in Chapter 6 we summarize our results and present our conclusion.}

